\begin{description}
\item[(1.1)] Look at the figure:

% FIGURE NEEDED: cluster sphere of radius R with Obs_1 at its centre and
% Obs_2 at distance D subtending aperture angle alpha (2012_th_sol.pdf,
% scanned solutions page S8)
\[
\tan\alpha \approx \alpha \approx \sin\alpha \approx \frac{2R}{D}
\]
\[
\Delta m = m_B - m_A = -2.5 \log\frac{F_B}{F_A}
\]
The brightness as seen by the astronomer ($F_A$) is the brightness of $N$
stars with absolute magnitude $M_0$ (and luminosity $L_0$) at a distance $D$
from him:
\[
F_A = N \frac{L_0}{4\pi D^2} \tag*{(4 pt)}
\]
We assume the cluster is homogeneous, so the density of stars $\rho$ inside
the cluster should be constant. The brightness coming from a spherical shell
with small thickness $\Delta R$ at a distance $R'$ is
\[
F_{\text{shell}}(R', \Delta R)
  = \rho V_{\text{shell}} \frac{L_0}{4\pi R'^2}
  = \frac{N}{\frac{4}{3}\pi R^3}\, 4\pi R'^2 \Delta R\, \frac{L_0}{4\pi R'^2}
  = \frac{3 N L_0 \Delta R}{4\pi R^3} \tag*{(5 pt)}
\]
The brightness doesn't depend on the radius of the shell $R'$. If we divide
the cluster in $n$ shells, $R = n \Delta R$, the brightness of the cluster as
seen by the astronomer 2 is: % sic: "astronomer 2" in the source; this is the
% flux seen by the observer at the centre (the biologist)
\[
F_B = n F_{\text{shell}} = n \frac{3 N L_0 \Delta R}{4\pi R^3}
    = \frac{3 N L_0}{4\pi R^2} \tag*{(4 pt)}
\]
So the difference in the magnitudes $\Delta m$ is:
\[
\Delta m = -2.5 \log\!\left(\frac{\;\frac{3NL_0}{4\pi R^2}\;}
                                 {\frac{NL_0}{4\pi D^2}}\right)
  = -2.5 \log\!\left[ 3 \left(\frac{D}{R}\right)^{\!2} \right]
  = 2.5 \log\!\left(\frac{\alpha^2}{12}\right) \tag*{(2 pt)}
\]

\item[(1.2)] The total energy received is proportional to the objective area:
\[
F'_A = F_A \frac{\pi \left(\frac{x}{2}\right)^2}
                {\pi \left(\frac{d_{\mathrm{pup}}}{2}\right)^2}
\quad\Rightarrow\quad
x = 6 \cdot 10^{-3} \sqrt{\frac{F'_A}{F_A}} \tag*{(3 pt)}
\]
\[
F'_A = F_B \quad\Rightarrow\quad
\frac{F'_A}{F_A} = \frac{F_B}{F_A} = \frac{12}{\alpha^2}
\quad\Rightarrow\quad
x = 6 \cdot 10^{-3} \sqrt{\frac{12}{\alpha^2}}
  = \frac{0.021}{\alpha}~m \tag*{(3 pt)}
\]

\item[(1.3)] In this case, a fraction
$\xi(d) = \dfrac{A(\alpha)}{4\pi d^2}$ of the light that comes from the
distance $d$ is seen by the biologist. $A(\alpha)$ is the area of the
interior of a cone with aperture $\alpha$ that intersects the sphere with
radius $d$ and center in the vertices of the cone (a spherical cap of radius
$d$ and height $D$): \hfill(2 pt)
\[
\cos\left(\frac{\alpha}{2}\right) = \frac{d - D}{d}
\quad\text{and}\quad
A(\alpha) = 2\pi D d
\;\Rightarrow\;
A(\alpha) = 2\pi d^2 \left[ 1 - \cos\left(\frac{\alpha}{2}\right) \right]
\;\Rightarrow\;
\xi = \sin^2\left(\frac{\alpha}{4}\right) \tag*{(4 pt)}
\]
So the fraction is independent of the distance $d$, and the fraction of light
that the biologist sees is $F'_B = \xi F_B$:
\[
\Delta m' = m'_B - m_A = -2.5 \log\frac{F'_B}{F_A}
  = -2.5 \log\!\left[ \frac{12 \sin^2\left(\frac{\alpha}{4}\right)}{\alpha^2} \right]
\tag*{(3 pt)}
\]
\end{description}
